\problemname{Strawberries}

\illustration{0.35}{jordgubbar.jpg}{}{}

Mikaela has harvested the first strawberries in all of Sweden, and wants to sell them for 
a high price. But in order to legally sell strawberries, she must first weigh them.
Unfortunately, Mikaela doesn't own a scale of any kind. The only tool she has that is up 
for this job is a vacuum cleaner.

There are $N$ strawberries, where the $i$th one has weight $w_i$. The strawberries are sorted by weight, 
so 
$$w_1 \leq w_2 \leq \dots \leq w_N.$$ 

All weights are integers between $1$ and $M$, where $M \leq N$. You can make queries of the form
\begin{itemize}
    \item \verb|? i s|: You set the strength of the vacuum cleaner to $s$, and use
    it on the strawberry with index $i$. If $w_i \leq s$, then the strawberry
    gets sucked up by the vacuum cleaner and is lost. Otherwise, the strawberry remains.
\end{itemize}

The goal is to find the weights of all of the strawberries that did not get lost, while also losing 
at most $M-1$ strawberries. In addition, in order to get full score you must use at most $300$ queries.

Note that you should not find the weights of the lost strawberries. Once a strawberry is lost, no more 
queries can be performed on it, but it still keeps its position in the list.

\section*{Interaction}

This is an interactive problem.
\begin{itemize}
    \item Your program should first read one line containing the integers $N$ and $M$ ($2 \leq M \leq N \leq 300$).
    \item Then, your program starts interacting with the grader. In order to make a query, print two integers $i$ and 
    $s$, on the format \verb|? i s|. These two integers must satisfy $1 \leq i \leq N$ and $1 \leq s \leq M$.

    After that, read one line containing a string, the answer to your query. This string is either \verb|borta| (lost)
    or \verb|kvar| (remaining), which indicates that strawberry $i$ was lost or remained. 

    If you try to print a query that doesn't follow the format, or try to query a strawberry that was already lost earlier,
    then you will get Wrong Answer.

    \item When you are ready to report the weights of all the remaining berries, print one line on the format \verb|! w1 w2 w3 ... wN|.
    The weights of lost berries should be printed as $-1$, otherwise you will get Wrong Answer.
\end{itemize}

If you use more than $10^5$ queries, you get Wrong Answer. If more than $M-1$ strawberries are lost, you get Wrong Answer. 
Reporting the final answer does not count as a query.

The grader is not adaptive. This means that the list of weights $w_1, w_2, \dots, w_N$ is fixed before the interaction starts.

\section*{Scoring}
Your solution will be tested on a set of test groups.
To earn points for a group, you must pass all test cases in that group.

For every test group, you will get $50\%$ of the score if your program uses between $301$ and 
$10^5$ queries. If you use at most $300$ queries, you get full score for that test group. 

\noindent
\begin{tabular}{| l | l | p{12cm} |}
  \hline
  \textbf{Group} & \textbf{Points} & \textbf{Constraints} \\ \hline
  $1$      & $16$      & $M = 2$  \\ \hline
  $2$      & $30$      & $M = N$ \\ \hline
  $3$      & $54$      & No additional constraints.  \\ \hline
\end{tabular}



\section*{Explanation of Sample}

In the first query, we use the vacuum cleaner on berry $2$, with $s = 1$. The berry gets lost. Then, we use the vacuum
cleaner on the last berry, with $s = 2$. The last berry remains.

Finally, we report the answer that $w_1 = 1$ and $w_3 = 3$. Note that we print $-1$ for the weight of berry $2$ 
because it got lost, even though we know that it must have had weight $1$.

Since $M = 3$ in this case, we can afford to lose at most $M-1 = 2$ berries. We only lost one berry, so this 
counts as Accepted.

